% !TeX root = ../../Book3.tex
% 纲要标题：### 18.5 正则点、切空间与光滑性
% 纲要位置：计划纲要.md，第 2078 行。

\section{正则点、切空间与光滑性}
\label{b3:sec:1805}

几何中的点由相应局部环描述，正则性比较该环的维数与一阶参数数目。相对光滑性还依赖基域；在不完美域上，即使环自身正则，改变基域后也可能出现幂零元。

% 纲要内容：
%
% * 正则点与切空间维数
% * 正则和基域上的光滑并非无条件相同
% * 非完美域上的区分例子
%
% ---
%

\subsection{正则点与切空间维数}
\label{b3:subsec:1805:01}

设 \(A\) 为域 \(k\) 上的有限型代数，\(\mathfrak p\in\operatorname{Spec}A\)。若 \(A_{\mathfrak p}\) 正则，就称 \(\mathfrak p\) 为\term[zhengzedian]{正则点}[regular point]。判定用的是该局部环的剩余域 \(\kappa(\mathfrak p)\)，即比较
\[
\dim A_{\mathfrak p}\quad\text{与}\quad
\dim_{\kappa(\mathfrak p)}
\mathfrak pA_{\mathfrak p}/(\mathfrak pA_{\mathfrak p})^2.
\]
在 \(k\) 有理闭点，第二个数也是 Zariski 切空间的维数，见\cref{b3:prop:tangent-cotangent-rational}。

\begin{proposition}\label{b3:prop:regular-hypersurface-rational-point}
设 \(k\) 为域，\(0\ne f\in k[X_1,\ldots,X_n]\)，且 \(a\in k^n\) 满足 \(f(a)=0\)。则超曲面 \(f=0\) 在 \(a\) 处的局部环正则，当且仅当至少一个 \(\partial f/\partial X_i(a)\ne0\)。
\end{proposition}
\begin{proof}
令 \(S=k[X_1,\ldots,X_n]_{(X_1-a_1,\ldots,X_n-a_n)}\)，其极大理想记为 \(\mathfrak n\)。\(S\) 正则且为整环，\(f\) 是非零因子，因此 \(S/(f)\) 的维数为 \(n-1\)。其极大理想模平方为
\(\mathfrak n/(\mathfrak n^2+(f))\)。\(f\) 在 \(\mathfrak n/\mathfrak n^2\) 中的像由一次 Taylor 项给出，系数恰为各偏导数在 \(a\) 处的值。若至少一项非零，商向量空间维数为 \(n-1\)；若全零，则为 \(n\)。和 Krull 维数比较即得结论。
\end{proof}

例如 \(y^2-x^3=0\) 在原点的偏导数全部为零，在任何特征都不正则。\(y-x^2=0\) 对 \(y\) 的偏导数为一，所以每个有理点正则。这一判断检测局部结构，并不依赖图形看起来是否有“尖角”。

\subsection{正则性与基域上的光滑性}
\label{b3:subsec:1805:02}

正则性是环自身的性质。光滑性则是基环到该环的同态的性质，\cref{b3:def:finite-presentation-infinitesimal-properties}用有限展示及平方零提升的存在性定义它。两者所问的问题不同，因此谈论光滑时必须注明基域或基环。

\ProseParagraphBreak
在前一命题的有理点处，如果某个偏导数不为零，可在一个主开邻域上使它可逆。对平方零提升问题，先任意提升各坐标，再用该偏导数的逆修正一个坐标，消掉方程值。这正是\secref{b3:sec:1404}的 Jacobian 线性修正，因而该主开邻域对 \(k\) 光滑。更一般的光滑性与几何正则性的等价判据在\chapref{b3:ch:23}建立，本章并不把那些尚待证明的结果用于正则环的同调或唯一分解证明。

\ProseParagraphBreak
已经得到的有理点判据不能直接照搬到剩余域不可分的点。\eqref{b3:eq:cotangent-residue-sequence}中的剩余域微分项此时可能非零，\(\Omega_{A/k}\) 的纤维不再只记录极大理想的一次部分。下面把这一差别转化为基域扩张后的局部环计算。

\subsection{非完美域上的区分例子}
\label{b3:subsec:1805:03}

取 \(k=\mathbb F_p(s)\)、\(K=k[u]/(u^p-s)\)。\cref{b3:exm:inseparable-tangent}已经说明，\(K\) 是域，\(\Omega_{K/k}\cong K\,\mathrm{d}u\ne0\)；\secref{b3:sec:1404}又以 \(k[T]/((T^p-s)^2)\) 中的平方零理想给出了不能提升的恒等映射。因此 \(K\) 是零维正则局部环，却不是光滑 \(k\) 代数。

\ProseParagraphBreak
现在把基域扩张到 \(K\)。有显式同构
\[
K\otimes_kK\cong K[T]/(T^p-s)
\cong K[v]/(v^p),\qquad v=T-u.
\]
最后一个环的唯一素理想为 \((v)\)，所以维数为零；而 \((v)/(v^2)\) 在 \(K\) 上维数为一，所以它不正则。这里环 \(K\) 自身完全没有零因子，基变换却产生幂零元。这说明“正则局部环”不包含相对于任意给定基域的可分性信息。

\ProseParagraphBreak
相反，\(k[x_1,\ldots,x_d]\) 扩张任何基域后仍是多项式环，坐标原点的局部环仍正则；它的平方零提升只需提升各坐标像。把这两类例子并列，可以看出环论正则性与相对光滑性在常见坐标例子中一致，却需要附加条件才可一般地等同。
